

\section{Le "Coup d'État" de 1880 : Jules Ferry et Michel Bréal imposent la "Version"}
space{0.3cm}
oindent

L'année 1880 ne marque pas seulement une étape législative dans l'histoire de la Troisième
République naissante ; elle constitue une véritable césure épistémologique et politique, un moment
de bascule violent que l'historiographie critique, à la suite des analyses de la section
I.1.\cite{I-1-2-ref2} de votre plan, qualifie à juste titre de « Coup d'État » pédagogique. Si la
mémoire collective retient volontiers les lois scolaires de 1881-1882 fondatrice de l'école
primaire, c'est pourtant dès 1880 que se joue, dans les hautes sphères de l'enseignement secondaire,
le destin de l'élite française. Sous la houlette de Jules Ferry, ministre de l'Instruction publique,
et de son inspirateur intellectuel, le linguiste Michel Bréal, une opération chirurgicale est menée
contre le cœur battant des humanités classiques traditionnelles : le « Discours latin » est aboli au
profit de la « Version » et de l'explication de texte.\cite{I-1-2-ref1}

Cette substitution, d'apparence technique, dissimule une ambition politique et sociale radicale : il
s'agit de « scientifiser » les humanités pour les arracher à l'influence de la rhétorique jésuite et
de l'éloquence d'Empire, jugées responsables de la décadence nationale et de la défaite de 1870 face
à la Prusse.\cite{I-1-2-ref3} En imposant la philologie et la traduction comme nouveaux critères de
l'excellence, Ferry et Bréal entendent former une nouvelle élite républicaine, caractérisée par
l'esprit critique, la rigueur positive et une maîtrise « nationale » de la langue française, en
rupture avec les « beaux parleurs » de l'Ancien Régime.\cite{I-1-2-ref5} Ce rapport se propose
d'explorer en détail les mécanismes, les acteurs et les conséquences de cette révolution culturelle
qui, en redéfinissant les exercices scolaires, a reconfiguré pour un siècle la distinction sociale
et intellectuelle en France.



\subsection{L'ancien régime pédagogique : l'hégémonie de la rhétorique et du « discours latin »}

Pour saisir la violence symbolique de la réforme de 1880, il convient d'abord d'analyser le système
qu'elle visait à détruire : un édifice pédagogique tricentenaire, hérité de la \textit{Ratio
Studiorum} des Jésuites et consolidé par l'Université napoléonienne, où la rhétorique régnait en
maître absolu.



\subsection{Le discours latin : la clé de voûte d'un système aristocratique}

Jusqu'à la fin des années 1870, l'enseignement secondaire français est structuré par une finalité
unique : la maîtrise de l'éloquence latine. L'exercice roi, celui qui sanctionne la fin des études
et détermine l'accès au baccalauréat, n'est ni la dissertation philosophique, ni l'analyse
littéraire, mais le « Discours latin » (ou composition latine). Cet exercice consistait à demander
aux élèves de se glisser dans la peau d'un personnage historique ou mythologique — par exemple, «
Alexandre le Grand haranguant ses troupes avant la bataille d'Arbèles » ou « Saint Louis s'adressant
à ses fils sur son lit de mort » — et de rédiger un discours en latin imitant le style de Cicéron,
de Tite-Live ou de Quinte-Curce.\cite{I-1-2-ref2}

Cet exercice reposait sur des principes pédagogiques et philosophiques précis, que les républicains
allaient bientôt dénoncer comme des vices :

\begin{itemize}\item \textbf{Le Principe d'Imitation (Mimesis) :} L'objectif de l'élève n'était pas de produire une pensée originale ou de découvrir une vérité historique, mais de reproduire le plus fidèlement possible des modèles canoniques (les \textit{auctores}). C'était une pédagogie de la répétition, où l'excellence se mesurait à la conformité au modèle.\cite{I-1-2-ref2}\item \textbf{La Primauté de la Forme sur le Fond :} Dans le discours latin, la vraisemblance rhétorique primait sur la vérité factuelle. Il s'agissait de trouver les arguments « probables » et les tournures élégantes (les \textit{colores}) pour persuader un auditoire fictif. L'élève apprenait à manipuler les émotions et les lieux communs (\textit{topos}) plutôt qu'à analyser des faits.\cite{I-1-2-ref5}\item \textbf{Une Culture de Caste :} La maîtrise de cet artifice demandait une imprégnation culturelle longue et coûteuse, accessible uniquement à la bourgeoisie aisée capable de financer des années d'humanités. Le discours latin fonctionnait comme un marqueur social, un signe de reconnaissance (« l'entre-soi ») distinguant l'honnête homme, capable de citer Virgile et de composer des vers, du vulgaire.\cite{I-1-2-ref8}\end{itemize}Comme le notait le philosophe Darlu en 1880, le discours latin était « la clé de voûte de l'ancien
édifice universitaire » : y toucher, c'était menacer l'équilibre de toute la formation des élites
françaises.\cite{I-1-2-ref7}



\subsection{L'ombre des jésuites et la critique politique}

Dans le contexte de lutte acharnée entre l'Église et la République naissante, le discours latin est
devenu, aux yeux des républicains, l'incarnation pédagogique du « péril clérical ». La Compagnie de
Jésus, principale force enseignante catholique concurrente de l'État, avait fait de la rhétorique
l'arme absolue de la Contre-Réforme pour séduire les esprits et les âmes.

\begin{itemize}\item \textbf{L'Arme de la Séduction :} Pour Jules Ferry, l'éloquence latine est suspecte car elle est l'outil par lequel les Jésuites forment des esprits dociles, habitués à défendre le faux comme le vrai, pourvu que la forme soit belle. C'est l'école de la sophistique et de l'hypocrisie.\cite{I-1-2-ref2}\item \textbf{Le Souvenir de l'Empire :} Cette méfiance est renforcée par le souvenir du Second Empire, régime perçu comme celui de la façade et du mensonge, qui avait lui aussi favorisé une culture de l'apparat. Après la chute de Napoléon III en 1870, la République, qui se veut le régime de la Raison, de la Science et de la Vertu, ne peut tolérer que ses élites soient formées à l'art du « verbiage ».\cite{I-1-2-ref11}\end{itemize}Le « Coup d'État » de 1880 n'est donc pas une simple réforme technique : c'est une opération de
désarmement politique. En supprimant le discours latin, Jules Ferry entend priver les congrégations
enseignantes de leur terrain d'excellence et imposer un nouveau terrain de bataille — la science et
la philologie — où l'Université laïque, armée par la recherche moderne, est certaine de
l'emporter.\cite{I-1-2-ref12}



\subsection{Le modèle allemand et la tentation scientifique : michel bréal, l'architecte intellectuel}

Si Jules Ferry a fourni la volonté politique et le bras séculier de l'État, c'est Michel Bréal qui a
fourni l'armature idéologique et épistémologique de la réforme. Figure centrale mais parfois
méconnue, Bréal est le véritable architecte de la « scientifisation » des humanités.



\subsection{Le traumatisme de 1870 et le complexe de sedan}

La réforme de 1880 est incompréhensible sans le traumatisme de la défaite de 1870. Une idée
obsédante parcourt l'élite intellectuelle française de l'après-guerre : « C'est le maître d'école
allemand qui a gagné à Sedan ».\cite{I-1-2-ref13} La France aurait été vaincue par la supériorité de
la science allemande, de son organisation, de sa rigueur, incarnées par l'Université de Berlin et
ses philologues.\cite{I-1-2-ref4}

Pour relever la France, il faut donc s'inspirer du vainqueur. Il faut importer en France
l'Altertumswissenschaft (la science de l'Antiquité), cette approche critique, historique et
philologique des textes, qui s'oppose à l'approche littéraire et rhétorique française.



\subsection{Michel bréal : le passeur de la philologie allemande}

Michel Bréal (1832-1915) est l'homme idéal pour opérer ce transfert culturel. Né à Landau (ville
bavaroise frontalière), bilingue, juif alsacien, il incarne cette frontière culturelle. Ancien élève
de Franz Bopp à Berlin, il est le traducteur en français de la monumentale Grammaire comparée des
langues indo-européennes.\cite{I-1-2-ref13}

Bréal introduit en France une idée révolutionnaire pour l'époque : la langue n'est pas seulement un
outil pour écrire (rhétorique), c'est un objet naturel et historique qu'il faut étudier
scientifiquement (linguistique).

Dans son ouvrage manifeste, Quelques mots sur l'instruction publique en France (1872), Bréal déploie
une critique dévastatrice de l'enseignement traditionnel. Il y consacre des chapitres entiers à
démolir les exercices rois : « Les vers latins », « Le discours latin » et « Le thème latin
».\cite{I-1-2-ref17}



\subsection{Le réquisitoire de bréal contre la rhétorique}

L'argumentation de Bréal, qui servira de bréviaire à Ferry, repose sur deux piliers : la vérité et
l'utilité.

\begin{itemize}\item \textbf{Contre l'Insincérité :} Bréal attaque le discours latin comme une école de mensonge. « On demande à l'élève d'exprimer des sentiments qu'il n'a pas, de parler au nom de personnages qu'il connaît à peine », écrit-il. Cela favorise le « verbiage », l'usage de mots vides et de formules toutes faites (les \textit{centons}). Pour Bréal, cet exercice corrompt le jugement et le sens de la vérité.\cite{I-1-2-ref19}\item \textbf{Pour la Méthode Scientifique :} À la place de la \textit{production} de textes latins factices, Bréal propose l'\textit{analyse} de textes authentiques. C'est l'avènement de l'explication de texte et de la \textbf{Version}. L'élève ne doit plus imiter Cicéron, il doit le \textit{comprendre}. Il doit disséquer la phrase, analyser les structures grammaticales, replacer le texte dans son contexte historique. C'est une démarche d'observation et de déduction, analogue à celle des sciences naturelles.\cite{I-1-2-ref21}\end{itemize}L'objectif est clair : transformer les humanités en « sciences de l'esprit ». Il s'agit de former
des esprits capables de critique, de distance et de précision, qualités jugées indispensables pour
l'élite d'une République moderne et industrielle.\cite{I-1-2-ref23}



\subsection{Le « coup d'état » de 1880 : mécanique et chronologie d'une révolution}

L'expression « Coup d'État » utilisée dans votre plan est particulièrement pertinente pour décrire
la méthode employée par Jules Ferry. Il ne s'est pas agi d'une évolution douce, mais d'une rupture
brutale, imposée par le haut, en contournant souvent les résistances corporatistes.



\subsection{Le contexte politique : la guerre scolaire}

En 1879-1880, les républicains prennent le contrôle total des institutions (Chambre, Sénat,
Présidence avec Jules Grévy). Jules Ferry, ministre de l'Instruction publique, lance simultanément
deux offensives :

1. \textbf{Contre les congrégations religieuses :} Les décrets de mars 1880 ordonnent la dissolution
des Jésuites et l'expulsion des congrégations non autorisées.

2. \textbf{Contre la pédagogie jésuite :} La réforme des programmes du lycée en 1880 est le volet
pédagogique de cette guerre. Il s'agit de laïciser non seulement le personnel, mais aussi l'esprit
même de l'enseignement.\cite{I-1-2-ref2}



\subsection{Les décrets de 1880 et la suppression du discours}

L'acte décisif est pris par le Conseil Supérieur de l'Instruction Publique (que Ferry vient de
réformer pour en exclure les évêques et y placer des universitaires acquis à sa cause, comme Bréal).

\begin{itemize}\item \textbf{Suppression du Discours Latin :} Les nouveaux programmes suppriment purement et simplement l'épreuve de composition latine au baccalauréat.\cite{I-1-2-ref25} C'est un choc immense pour le corps professoral traditionnel. Du jour au lendemain, ce qui occupait 10 à 15 heures de cours par semaine devient inutile pour l'examen.\item \textbf{Suppression des Vers Latins :} L'exercice poétique, jugé frivole et aristocratique, est également éliminé.\cite{I-1-2-ref2}\item \textbf{Montée en Puissance de la Version et de la Composition Française :} Pour remplacer le discours, Ferry instaure deux nouvelles épreuves reines :\item La \textbf{Version latine} (traduction du latin vers le français), qui devient l'épreuve de discrimination par excellence.\item La \textbf{Composition française} (ancêtre de la dissertation), qui remplace le discours latin comme exercice de rhétorique, mais une rhétorique désormais nationale et laïque.\cite{I-1-2-ref5}\end{itemize}

\subsection{La résistance : « la question du latin »}

Cette réforme suscite une levée de boucliers, connue sous le nom de « Question du Latin ». Des
intellectuels conservateurs, mais aussi certains libéraux, accusent Ferry de « tuer les humanités »
et de rabaisser le niveau culturel de la France.

\begin{itemize}\item En 1885, Raoul Frary publie \textit{La Question du latin}, un pamphlet qui attaque violemment le fétichisme du latin, soutenant paradoxalement la réforme de Ferry tout en allant plus loin (il prône la suppression totale du latin obligatoire).\cite{I-1-2-ref28}\item Les défenseurs de la tradition (comme Ferdinand Brunetière) voient dans la suppression du discours latin une « américanisation » ou une « germanisation » de l'esprit français, craignant que l'élite ne devienne une caste de techniciens sans âme.\cite{I-1-2-ref2}\end{itemize}

\subsection{La version et la formation de l'élite républicaine : distinction et nationalisme}

Pourquoi Jules Ferry et Michel Bréal ont-ils choisi la \textbf{Version} comme nouvel instrument de
sélection? Comment cet exercice de traduction a-t-il servi à former l'élite républicaine spécifique
de la IIIe République?



\subsection{La scientifisation de la sélection : l'épreuve de vérité}

Contrairement au discours latin, où l'on pouvait briller par le style tout en faisant des contresens
historiques, la Version est impitoyable.

\begin{itemize}\item \textbf{Rigueur et Exactitude :} La Version exige une compréhension totale de la grammaire et du lexique. C'est une épreuve binaire : c'est juste ou c'est faux. Elle introduit une forme de « justice » méritocratique brutale, compatible avec l'idéal républicain d'égalité devant le concours.\cite{I-1-2-ref22}\item \textbf{L'Esprit d'Analyse :} Bréal présente la Version comme une gymnastique intellectuelle supérieure. Face à une phrase de Tacite ou de Cicéron, l'élève doit émettre des hypothèses, vérifier les cas, analyser les relations syntaxiques. C'est une formation à l'esprit d'enquête et de déduction. Le futur haut fonctionnaire ou ingénieur formé par la Version aura appris à ne pas se payer de mots, mais à chercher le sens exact des textes et des lois.\cite{I-1-2-ref21}\end{itemize}

\subsection{Le nationalisme linguistique : apprendre le français par le latin}

Un paradoxe fondamental de cette réforme est que le latin sert désormais, avant tout, à apprendre le
français. C'est la thèse développée par André Chervel et Renée Balibar.

\begin{itemize}\item \textbf{La « Traduction » vers le National :} Dans l'exercice de Version, l'effort principal n'est pas tant de comprendre le latin que de trouver l'équivalent exact en français. Cela oblige l'élève à explorer les ressources de sa propre langue maternelle.\item \textbf{Unification Linguistique :} À une époque où les patois sont encore vivaces, la Version impose un français normé, classique, académique. Elle participe à la construction de la langue nationale, outil indispensable de l'unité républicaine. Le latin de Ferry est un détour pédagogique pour fabriquer des citoyens français possédant une maîtrise parfaite de la langue de l'État.\cite{I-1-2-ref1}\end{itemize}

\subsection{La distinction sociale : l'analyse de pierre bourdieu}

Si la réforme se veut méritocratique, elle n'en demeure pas moins un puissant outil de reproduction
sociale, comme l'a analysé Pierre Bourdieu.

\begin{itemize}\item \textbf{La Barrière de l'Abstraction :} En remplaçant le capital culturel « mondain » (le style, l'éloquence) par un capital culturel « scolaire » (la rigueur grammaticale, la philologie), la réforme déplace la barrière de classe mais ne l'abolit pas. La réussite en Version latine nécessite un habitus de rigueur, de calme et une longue fréquentation des textes écrits, apanage des classes bourgeoises.\cite{I-1-2-ref33}\item \textbf{La Noblesse d'État :} La Version devient l'épreuve reine du Concours Général et de l'entrée à l'École Normale Supérieure (la rue d'Ulm). Elle sélectionne une « noblesse d'État » fondée sur la compétence scolaire et le diplôme, remplaçant l'ancienne élite des notables. Bourdieu lui-même, pur produit de cette méritocratie (lauréat du Concours Général en version latine), incarne cette nouvelle élite républicaine sélectionnée par la traduction.\cite{I-1-2-ref35}\end{itemize}

\subsection{Synthèse des données et tableaux comparatifs}

Pour visualiser l'ampleur de la rupture de 1880, nous proposons ci-dessous une comparaison
structurelle des deux systèmes éducatifs.



\subsection{Tableau 1 : comparaison du modèle rhétorique (avant 1880\) et du modèle philologique (après 1880\)}

\begin{table}[h!]\small\centering\begin{tabular}{| p{2.3cm} | p{3.5cm} | p{3.5cm} |}\hline\textbf{Dimension} & \textbf{Ancien Régime Pédagogique (Jésuites / Empire)} & \textbf{Nouveau Régime Républicain (Ferry / Bréal)} \\ \hline\textbf{Exercice Roi} & Le Discours Latin (Composition) \& Vers Latins & La Version Latine (Traduction) \& Explication de texte \\ \hline\textbf{Objectif Pédagogique} & \textbf{Mimesis} (Imitation des Anciens) & \textbf{Analyse} (Compréhension critique) \\ \hline\textbf{Qualité Valorisée} & L'Éloquence, le Style, l'Amplification & La Rigueur, l'Exactitude, la Vérité historique \\ \hline\textbf{Relation au Texte} & Le texte est un \textbf{modèle} à reproduire & Le texte est un \textbf{objet} à étudier scientifiquement \\ \hline\textbf{Modèle Étranger} & Rome (Humanisme latin) & Allemagne (Philologie et \textit{Altertumswissenschaft}) \\ \hline\textbf{Idéal Social} & L'Orateur (Avocat, Tribun, Prêtre) & Le Savant (Professeur, Ingénieur, Fonctionnaire) \\ \hline\textbf{Fonction du Latin} & Langue de production (écrire \textit{en} latin) & Langue de culture (lire \textit{du} latin pour former le français) \\ \hline\end{tabular}\end{table}

\subsection{Tableau 2 : les acteurs clés du « coup d'état » de 1880}

\begin{table}[h!]\small\centering\begin{tabular}{| p{2cm} | p{2.2cm} | p{2.2cm} | p{2.2cm} |}\hline\textbf{Acteur} & \textbf{Rôle et Fonction} & \textbf{Apport à la Réforme} & \textbf{Source Clé} \\ \hline\textbf{Jules Ferry} & Ministre de l'Instruction Publique & Volonté politique, anticlericalisme, imposition des décrets. & 2 \\ \hline\textbf{Michel Bréal} & Linguiste, Professeur au Collège de France & Architecte intellectuel, importateur du modèle allemand, critique du \og verbiage\fg{}. & 4 \\ \hline\textbf{Les Jésuites} & Congrégation enseignante (opposants) & Cible de la réforme, défenseurs de la rhétorique et du discours latin. & 2 \\ \hline\textbf{Raoul Frary} & Journaliste et essayiste & Critique radical, auteur de \textit{La Question du latin} (1885), pousse la logique de Ferry. & 28 \\ \hline\textbf{Pierre Bourdieu} & Sociologue (XXe s.) & Analyste \textit{a posteriori} de la fonction sociale de sélection de la Version. & 33 \\ \hline\end{tabular}\end{table}

\subsection{Conséquences à long terme et débats historiographiques}

La réussite du « Coup d'État » de 1880 ne doit pas masquer les résistances et les adaptations qui
ont suivi. L'historiographie contemporaine permet de nuancer le triomphe de la philologie.



\subsection{La persistance du « signe » (françoise waquet)}

Comme le démontre Françoise Waquet dans \textit{Le latin ou l'empire d'un signe}, le latin n'a pas
disparu, il s'est métamorphosé. Si la prétention scientifique de Bréal a légitimé le maintien du
latin, dans la pratique sociale, le latin est resté un « signe » de classe. Les familles bourgeoises
ont continué à plébisciter la section classique non par amour de la philologie, mais parce qu'elle
garantissait l'entre-soi social. La « Version » a simplement remplacé le « Discours » comme brevet
de bourgeoisie.\cite{I-1-2-ref38}



\subsection{La création d'une discipline scolaire autonome (andré chervel)}

André Chervel, historien des disciplines scolaires, souligne que l'école a digéré la réforme à sa
manière. La « Version » scolaire n'est pas devenue la philologie universitaire allemande. Elle est
devenue un exercice sui generis, avec ses propres règles (la « belle traduction »), ses propres
codes (les « faux sens », « contresens », « barbarismes »). L'école a créé une « culture scolaire »
qui a fini par s'autonomiser de la science universitaire qu'elle était censée imiter. Le latin
scolaire de 1900 ou 1950 n'est plus celui des Jésuites, mais il n'est pas non plus tout à fait la
linguistique de Bopp : c'est le latin de la « classe de rhétorique » républicaine.\cite{I-1-2-ref27}



\subsection{L'esprit « khâgneux » et la dissertation}

Enfin, l'héritage le plus durable de 1880 est sans doute la consécration de la \textbf{Dissertation}
(composition française) et de la \textbf{Version} comme les deux jambes de l'excellence française
(prépa littéraire, Khâgne). Ce modèle, fondé sur la culture générale, l'aptitude à synthétiser et à
traduire, a façonné l'élite administrative et politique française (ENA, Sciences Po) bien au-delà de
la IIIe République. En ce sens, Jules Ferry et Michel Bréal ont réussi leur pari : ils ont créé un
moule intellectuel qui a survécu aux régimes successifs.\cite{I-1-2-ref25}

En conclusion, la section I.1.\cite{I-1-2-ref2} de votre plan, identifiant le moment 1880 comme un «
Coup d'État », touche au cœur de la fabrique de l'histoire française. Ce n'était pas seulement une
réforme de programmes ; c'était une redéfinition de ce que signifie « penser » et « être cultivé »
en France.

En remplaçant l'éloquence (la parole) par la version (l'analyse), Ferry et Bréal ont opéré une
\textbf{rupture anthropologique}. Ils ont substitué à l'homme de la tribune (l'idéal révolutionnaire
et romantique) l'homme du bureau et du laboratoire (l'idéal positiviste et républicain). Ils ont
armé la République d'une élite formée à la rigueur, protégée du cléricalisme par le bouclier de la
science critique, et unifiée par une langue nationale épurée par le passage obligé de la traduction.

Si la « scientifisation » totale des humanités est restée en partie un mythe — le littéraire et
l'esthétique résistant toujours —, l'institutionnalisation de la Version a bel et bien créé cette «
élite républicaine » qui allait diriger la France jusqu'à la fin du XXe siècle. Jules Ferry ne s'est
pas contenté de construire des écoles ; avec Michel Bréal, il a construit l'esprit même de ceux qui
allaient les diriger. La Version latine fut leur outil, et le Lycée, leur laboratoire.

